% Part: first-order-logic
% Chapter: sequent-calculus
% Section: proof-theoretic-notions

\documentclass[../../../include/open-logic-section]{subfiles}

\begin{document}

\begin{editorial}
  本节汇集了自然演绎中可证性关系与一致性的定义。
\end{editorial}

\iftag{FOL}
      {\olfileid{fol}{seq}{ptn}}
      {\olfileid{pl}{seq}{ptn}}

\olsection{证明论概念}

\begin{explain}
正如我们曾定义过若干重要的语义概念（有效性、语义后承、可满足性）一样，
我们现在来定义相应的\emph{证明论概念}。
这些概念不是通过语句在结构中是否被满足来定义的，而是通过某些矢列的可推导性或不可推导性来定义的。
一个重要的发现是，这些概念与语义概念恰好重合。
这正是\emph{可靠性定理}与\emph{完备性定理}所揭示的内容。
\end{explain}

\begin{defn}[定理]
语句~$!A$ 是\emph{定理}，当且仅当在~$\Log{LK}$ 中存在矢列 $\quad \Sequent !A$ 的一个推导。
若 $!A$ 是定理，则记作 $\Proves !A$；若不是，则记作 $\Proves/ !A$。
\end{defn}

\begin{defn}[可推导性]
语句~$!A$ 可从语句集~$\Gamma$ \emph{推导}（记作 $\Gamma \Proves !A$），
当且仅当存在一个有限子集~$\Gamma_0 \subseteq \Gamma$ 以及由~$\Gamma_0$ 中语句构成的一个序列 $\Gamma_0'$，
使得 $\Log{LK}$ 可推导出矢列 $\Gamma_0' \Sequent !A$。
若 $!A$ 不能从 $\Gamma$ 推导，则记作 $\Gamma \Proves/ !A$。
\end{defn}

由于收缩、弱化与交换规则的存在，$\Gamma_0'$ 中语句的顺序和数量并不重要：如果矢列 $\Gamma_0' \Sequent !A$ 是可推导的，那么对于任何与~$\Gamma_0'$ 包含相同语句的 $\Gamma_0''$，$\Gamma_0'' \Sequent !A$ 也是可推导的。例如，如果 $\Gamma_0 = \{!B, !C\}$
那么 $\Gamma_0' = \tuple{!B, !B, !C}$ 和 $\Gamma_0'' =
\tuple{!C, !C, !B}$ 都是仅包含 $\Gamma_0$ 中语句的序列。如果以其中之一为前件的矢列是可推导的，那么以另一个为前件的矢列也是可推导的，例如：

\begin{prooftree}
  \AxiomC{}
  \Deduce$!B, !B, !C \fCenter !A$
  \RightLabel{\LeftR{\Contraction}}
  \UnaryInf$!B, !C \fCenter !A$
  \RightLabel{\LeftR{\Exchange}}
  \UnaryInf$!C, !B \fCenter !A$
  \RightLabel{\LeftR{\Weakening}}
  \UnaryInf$!C, !C, !B \fCenter !A$
\end{prooftree}

从现在起，若 $\Gamma_0$ 是一个有限语句集，我们将用 $\Gamma_0 \Sequent !A$ 表示任何前件为~$\Gamma_0$ 中语句序列的矢列，并在必要时默认隐含了收缩、交换和弱化步骤。

\begin{defn}[一致性]
语句集~$\Gamma$ 是\emph{不一致的}，当且仅当存在一个有限子集~$\Gamma_0 \subseteq \Gamma$，使得 $\Log{LK}$ 可推导出矢列 $\Gamma_0 \Sequent \quad$。
如果 $\Gamma$ 不是不一致的，即对于每个有限子集 $\Gamma_0 \subseteq \Gamma$，$\Log{LK}$ 都不能推导出 $\Gamma_0 \Sequent \quad$，则称它是\emph{一致的}。
\end{defn}

\begin{prop}[自反性]
\ollabel{prop:reflexivity}
如果 $!A \in \Gamma$，那么 $\Gamma \Proves !A$。
\end{prop}

\begin{proof}
初始矢列 $!A \Sequent !A$ 是可推导的，且 $\{!A\} \subseteq \Gamma$。
\end{proof}

\begin{prop}[单调性]
\ollabel{prop:monotonicity}
如果 $\Gamma \subseteq \Delta$ 且 $\Gamma \Proves !A$，那么 $\Delta \Proves !A$。
\end{prop}

\begin{proof}
假设 $\Gamma \Proves !A$，即存在有限子集 $\Gamma_0 \subseteq \Gamma$ 使得矢列 $\Gamma_0 \Sequent !A$ 是可推导的。由于 $\Gamma \subseteq \Delta$，$\Gamma_0$ 也是 $\Delta$ 的一个有限子集。因此，$\Gamma_0 \Sequent !A$ 的推导同样证明了 $\Delta \Proves !A$。
\end{proof}

\begin{prop}[传递性]
\ollabel{prop:transitivity}
如果 $\Gamma \Proves !A$ 且 $\{!A\} \cup \Delta \Proves
!B$，那么 $\Gamma \cup \Delta \Proves !B$。
\end{prop}

\begin{proof}
如果 $\Gamma \Proves !A$，则存在有限子集 $\Gamma_0 \subseteq \Gamma$ 以及 $\Gamma_0 \Sequent !A$ 的一个推导~$\pi_0$。
如果 $\{!A\}
\cup \Delta \Proves !B$，则对于某个有限子集 $\Delta_0 \subseteq \Delta$，存在 $!A, \Delta_0 \Sequent !B$ 的一个推导~$\pi_1$。
考虑如下推导：
\begin{prooftree}
  \AxiomC{}
  \RightLabel{$\pi_0$}
  \Deduce$\Gamma_0 \fCenter !A$
  \AxiomC{}
  \RightLabel{$\pi_1$}
  \Deduce$!A, \Delta_0 \fCenter !B$
  \RightLabel{\Cut}
  \BinaryInf$\Gamma_0, \Delta_0 \fCenter !B$
\end{prooftree}
由于 $\Gamma_0 \cup \Delta_0 \subseteq \Gamma \cup \Delta$，这就证明了 $\Gamma \cup \Delta \Proves !B$。
\end{proof}

注意，这特别地意味着：如果 $\Gamma \Proves !A$ 且 $!A \Proves !B$，那么 $\Gamma \Proves !B$。
由此也可得，如果 $!A_1,
\dots, !A_n \Proves !B$ 且对每个~$i$ 有 $\Gamma \Proves !A_i$，那么 $\Gamma \Proves !B$。

\begin{prop}
\ollabel{prop:incons}
$\Gamma$ 是不一致的，当且仅当 $\Gamma$ 可推导出任一语句~$!A$。
\end{prop}

\begin{proof}
习题。
\end{proof}

\begin{prob}
证明 \olref[fol][seq][ptn]{prop:incons}。
\end{prob}

\begin{prop}[紧致性]
  \ollabel{prop:proves-compact}
  \begin{enumerate}
  \item 如果 $\Gamma \Proves !A$，则存在 $\Gamma$ 的一个有限子集 $\Gamma_0 \subseteq \Gamma$，使得 $\Gamma_0 \Proves !A$。
  \item 如果 $\Gamma$ 的每个有限子集都是一致的，那么 $\Gamma$ 就是一致的。
  \end{enumerate}
\end{prop}

\begin{proof}
  \begin{enumerate}
    \item 如果 $\Gamma \Proves !A$，则存在 $\Gamma$ 的一个有限子集 $\Gamma_0 \subseteq \Gamma$，使得矢列 $\Gamma_0 \Sequent !A$ 有一个推导。因此，$\Gamma_0 \Proves !A$。
    \item 如果 $\Gamma$ 是不一致的，则存在一个有限子集 $\Gamma_0 \subseteq \Gamma$，使得 $\Log{LK}$ 能推导出 $\Gamma_0 \Sequent \quad$。这样一来，$\Gamma_0$ 就是 $\Gamma$ 的一个不一致的有限子集。
  \end{enumerate}
\end{proof}

\end{document}